\documentclass[11pt]{article}
\usepackage[margin=0.88in]{geometry}
\usepackage{amsmath,amssymb,lmodern}
\usepackage[T1]{fontenc}
\usepackage{needspace}
\usepackage[colorlinks=true,linkcolor=blue,urlcolor=blue]{hyperref}
\setlength{\emergencystretch}{3em}
\setlength{\parskip}{0.35em}
\setcounter{secnumdepth}{0}
\providecommand{\tightlist}{\setlength{\itemsep}{0pt}\setlength{\parskip}{0pt}}
\title{Positive radial series and sharp logarithmic concavity on balls}
\author{Henry Zweiman}
\date{September 15, 2026}
\begin{document}
\maketitle
\Needspace{8\baselineskip}\section{Abstract}\label{abstract}

We study positive Dirichlet solutions of \(-\Delta u=u\log u^2\) on a
ball in any dimension. We prove that \(-\log(u(r)/u(0))\) has a power
series in \(r^2\) with strictly positive coefficients and convergence
radius equal to the square of the ball radius. The coefficients follow
an explicit positive recurrence and are polynomials with nonnegative
coefficients in the excess of the central logarithmic height above the
Gaussian threshold. This structure proves concavity of
\(-\sqrt{-\log(u/\|u\|_\infty)}\), answering a question of Gallo,
Mosconi and Squassina. The exponent \(1/2\) is sharp among positive
powers of the negative normalized logarithm. We also obtain a
quantitative lower bound for its Hessian before taking the square root,
positive power concavity of \(u\) on every ball, and a comparison of
normalized profiles. The argument gives an alternative proof of the
known shooting order and ball uniqueness theorem. The essential step is
to pass from a local positive recurrence to a convergent representation
throughout the domain.

\textbf{Mathematics Subject Classification (2020).} 35B45, 35J61, 34B15,
26A51.

\textbf{Keywords.} Logarithmic Schrödinger equation; concavity; radial
solutions; positive power series; shooting method.

\Needspace{8\baselineskip}\section{1. Introduction and main
results}\label{introduction-and-main-results}

Let \(N\geq1\), \(R>0\), and \(B_R=\{x\in\mathbb R^N:|x|<R\}\). We
consider

\[
-\Delta u=u\log u^2\quad\text{in }B_R,
\qquad u>0\quad\text{in }B_R,
\qquad u=0\quad\text{on }\partial B_R,
\tag{1.1}
\]

where \(u\in C^2(B_R)\cap C(\overline{B_R})\). All concavity statements
about a logarithmic transform concern the open ball. Statements about
\(u^\alpha\) extend to the closed ball by continuity.

Gallo, Mosconi and Squassina {[}GMS{]} establish existence of a solution
with locally negative-definite logarithmic Hessian on every bounded
convex domain. Their Theorem 6.5 gives radial symmetry and radial
decrease of every positive solution on a ball. In Sections 6.2-\/-6.3
they ask for stronger concavity in balls: positive power concavity of
\(u\), and concavity of the negative square root of its negative
normalized logarithm. Their Theorem 6.6 proves both properties in
dimension one. The second question remains explicitly stated on page 44
of the final article.

Our main result is a positive-series representation that answers both
questions in every dimension. Write \(u(x)=U(r)\), \(r=|x|\), and set

\[
M=\log U(0),\qquad
w(r)=-\log\frac{U(r)}{U(0)},\qquad
A=\frac{2M}{N}.
\tag{1.2}
\]

\Needspace{6\baselineskip}\textbf{Theorem 1.1 (Positive radial series).} Every solution of (1.1)
satisfies \(A>1\). There are strictly positive coefficients \(c_n\),
\(n\geq0\), such that

\[
\frac{w'(r)}r=\sum_{n=0}^\infty c_n r^{2n},
\qquad
w(r)=\sum_{j=1}^\infty\frac{c_{j-1}}{2j}r^{2j},
\qquad 0\leq r<R.
\tag{1.3}
\]

The quotient at \(r=0\) is interpreted by continuity. Both series,
regarded as series in \(r^2\), have radius of convergence exactly
\(R^2\). Their coefficients are determined by

\[
c_0=A,\qquad
c_{n+1}=
\frac{(n+1)\displaystyle\sum_{i=0}^n c_i c_{n-i}-c_n}
{(n+1)(2n+N+2)}.
\tag{1.4}
\]

In particular, \(Z(t)=w'(\sqrt t)/\sqrt t\) is absolutely monotone on
\([0,R^2)\): every derivative is strictly positive there. The same is
true of every derivative of positive order of \(W(t)=w(\sqrt t)\). In
addition,

\[
D^2w(|x|)\succeq A I_N,
\tag{1.5}
\]

with strict inequality as quadratic forms when \(x\neq0\). Here
\(\succeq\) denotes the order on symmetric matrices.

\Needspace{6\baselineskip}\textbf{Theorem 1.2 (Sharp transformed concavity).} For every solution
of (1.1),

\[
x\longmapsto -\left(-\log\frac{u(x)}{\|u\|_\infty}\right)^\beta
\quad\text{is concave on }B_R
\quad\Longleftrightarrow\quad \beta\geq\frac12,
\qquad \beta>0.
\tag{1.6}
\]

For \(\beta=1/2\) its Hessian is negative definite away from the center.
At the center it is continuous and has a norm-type cusp; no
differentiability assertion there is intended.

\Needspace{6\baselineskip}\textbf{Theorem 1.3 (Positive power concavity).} There exists a number
\(\alpha_*(U)\in(0,1)\) such that, for \(\alpha>0\),

\[
u^\alpha\text{ is concave on }\overline{B_R}
\quad\Longleftrightarrow\quad
\alpha\leq\alpha_*(U).
\tag{1.7}
\]

It is characterized by an attained interior minimum,

\[
\alpha_*(U)=\min_{0<r<R}\frac{w''(r)}{w'(r)^2}.
\tag{1.8}
\]

The formula is implicit. We do not assert a positive exponent uniform
over all radii, or monotonicity of this optimal exponent as a function
of the radius.

\subsection{1.1. Prior results and the scope of the
contribution}\label{prior-results-and-the-scope-of-the-contribution}

The amplitude threshold, uniqueness in balls, and strict ordering of
first zeros are prior results. For \(k=1\) and \(\beta=0\), Ben Chrouda
{[}BC, Theorem 1.2 and Corollaries 1.3 and 1.5{]} proves these facts for
\(-\Delta v=v\log v\). The change \(v(y)=u(y/\sqrt2)\) gives the
normalization (1.1). The equality case in the amplitude threshold is
excluded by the Gaussian solution and uniqueness of the central
initial-value problem. We give new proofs of these particular
consequences because they fit the recurrence argument, but do not claim
their first discovery.

Liu, Sun and Zou {[}LSZ{]} prove a broader shooting classification of
whole-space nodal solutions; in particular, their argument establishes
strict motion of zeros with the initial value. Gallo, Moraschi and
Squassina {[}GMoS{]} develop parabolic concavity principles whose
logarithmic example yields preservation of logarithmic concavity.
Neither of these results supplies the positive radial series used here.
Our priority claim concerns that structure and the resulting stronger
concavity, rather than existence, radial symmetry, Gaussian ground-state
classification, or shooting uniqueness. Current versions of these
sources were compared on September 15, 2026.

The proof of Theorem 1.1 has three stages. An integrating factor puts
the logarithmic radial derivative above its Gaussian value. Its equation
then produces positive Taylor coefficients. Finally, positivity and a
real-axis singularity argument show that the series represents the
solution all the way to its first zero. Convexity of a Hilbert-space
norm proves Theorem 1.2. A separate boundary argument, which
accommodates the singular derivative of \(u\log u^2\) at zero, completes
Theorem 1.3.

\Needspace{8\baselineskip}\section{2. The radial equation and the Gaussian
threshold}\label{the-radial-equation-and-the-gaussian-threshold}

We use radial symmetry and radial decrease from {[}GMS, Theorem 6.5{]}.
These are nontrivial inputs: the derivative of the nonlinearity is
unbounded below near zero, so an unqualified appeal to the usual locally
Lipschitz moving-plane theorem would be inappropriate. Existence follows
from {[}GMS, Theorem 1.1{]}, with its one-dimensional result covering
\(N=1\). The rest of the argument applies to every positive radial
decreasing solution.

Put \(f(s)=s\log s^2\) for \(s>0\), and \(f(0)=0\). The radial equation
and its integral form are

\[
U''+\frac{N-1}{r}U'+f(U)=0,
\qquad
U'(r)=-r^{1-N}\int_0^r s^{N-1}f(U(s))\,ds.
\tag{2.1}
\]

The central initial-value problem is unique as long as its solution
stays positive. Indeed, integrating (2.1) once more gives a Volterra
equation. On a sufficiently small interval the Lipschitz bound for \(f\)
near a prescribed positive central value makes its integral operator a
contraction, since the norm of the double integral is bounded by a
constant times \(r^2\). Ordinary uniqueness then continues the solution
away from zero. The same identity, using smoothness of \(f\) near a
positive central value, gives the usual even expansions at the center.

Since \(U\) has its maximum at zero, (2.1) implies \(f(U(0))\geq0\),
hence \(M\geq0\). If \(M=0\), central uniqueness would give
\(U\equiv1\), which violates the boundary condition. Thus \(M>0\).

In fact \(U'(r)<0\) for \(0<r<R\). If \(U'\) vanished there, its
nonpositivity would imply \(U''=0\) at that point. Equation (2.1) would
give \(U=1\), and ordinary uniqueness with data \((U,U')=(1,0)\) would
again force the constant solution.

In terms of (1.2), equation (2.1) becomes

\[
w''+\frac{N-1}{r}w'-(w')^2=2M-2w.
\tag{2.2}
\]

Define \(z(r)=w'(r)/r\) for \(r>0\) and extend it to the center. The
central expansion gives

\[
z(0)=\frac{2M}{N}=A,\qquad z'(r)=O(r).
\tag{2.3}
\]

Differentiation of (2.2) yields

\[
z''+\left(\frac{N+1}{r}-2rz\right)z'=2z(z-1).
\tag{2.4}
\]

Equivalently,

\[
\left(r^{N+1}e^{-2w}z'\right)'
=2r^{N+1}e^{-2w}z(z-1).
\tag{2.5}
\]

\Needspace{6\baselineskip}\textbf{Lemma 2.1.} We have \(A>1\) and \(z'(r)>0\) for all \(0<r<R\).

\emph{Proof.} Suppose first that \(0<A<1\). Integrating (2.5) from zero
shows that \(z'<0\) while \(0<z<1\). The boundary term at zero vanishes
by (2.3). The function cannot first reach \(1\), since it is decreasing.
It cannot first reach zero either: the integral in (2.5) would then give
a strictly negative derivative, whereas \(z=w'/r>0\) throughout the
interior. Thus \(0<z\leq A\) on \([0,R)\), giving \(w(r)\leq Ar^2/2\).
This contradicts \(w(r)\to\infty\) as \(r\uparrow R\).

If \(A=1\), the central solution is

\[
U(r)=\exp\left(\frac N2-\frac{r^2}2\right),
\tag{2.6}
\]

by the uniqueness following (2.1). It has no finite zero. Hence \(A>1\).
Now (2.5) gives \(z'>0\) as long as \(z>1\), and this strict increase
prevents any first exit from that region. \(\square\)

The eigenvalues of \(D^2w(|x|)\) are \(w''=z+rz'\) in the radial
direction and \(z\) in the tangential directions. Lemma 2.1 proves
(1.5), including its strict form away from the center. In dimension one
only the radial eigenvalue is present.

\Needspace{8\baselineskip}\section{3. A positive recurrence valid throughout the
ball}\label{a-positive-recurrence-valid-throughout-the-ball}

Let \(t=r^2\) and \(Z(t)=z(\sqrt t)\). Equation (2.4) becomes

\[
2tZ''+(N+2)Z'-2tZZ'=Z(Z-1).
\tag{3.1}
\]

\Needspace{6\baselineskip}\textbf{Lemma 3.1 (Local construction).} For each \(A>1\), recurrence
(1.4) defines strictly positive coefficients satisfying

\[
0<c_n\leq A^{n+1},\qquad n\geq0.
\tag{3.2}
\]

The resulting series solves (3.1) near zero and agrees with the
transformed radial solution.

\emph{Proof.} The coefficient of \(t^n\) in the linear derivative terms
of (3.1) is \((n+1)(2n+N+2)c_{n+1}\). In the nonlinear terms, symmetry
of convolution gives

\[
[t^n](Z^2+2tZZ')=(n+1)\sum_{i=0}^n c_i c_{n-i},
\tag{3.3}
\]

which proves (1.4). In particular \(c_1=A(A-1)/(N+2)>0\). For
\(n\geq1\), its numerator can be written

\[
\bigl(2(n+1)A-1\bigr)c_n
+(n+1)\sum_{i=1}^{n-1}c_i c_{n-i},
\tag{3.4}
\]

so positivity follows by induction. Dropping the negative term in (1.4)
and using the inductive upper bounds gives

\[
c_{n+1}\leq\frac{(n+1)A^{n+2}}{2n+N+2}\leq A^{n+2}.
\tag{3.5}
\]

Consequently the series converges for \(|t|<1/A\), and termwise
differentiation verifies (3.1).

Integrate \(w'(r)=rZ(r^2)\) with \(w(0)=0\). The derivative of the
difference between the two sides of (2.2) vanishes by (3.1). Its
limiting value at zero is \(NA-2M=0\). Thus (2.2) holds. Setting
\(U(r)=e^{M-w(r)}\) recovers the radial equation with the prescribed
central data. The uniqueness argument after (2.1) identifies this
constructed solution with the given one. \(\square\)

We next state the positive-series continuation principle, including its
short proof.

\Needspace{6\baselineskip}\textbf{Lemma 3.2 (Positive real singularity).} Suppose
\(F(t)=\sum_{n\geq0}a_nt^n\) has \(a_n\geq0\) and finite positive radius
of convergence \(\rho\). It cannot admit a holomorphic continuation to a
neighborhood of the positive real point \(\rho\).

\emph{Proof.} If such a continuation existed, continuity of its
derivatives and monotone convergence as \(t\uparrow\rho\) would give,
for every \(k\geq0\),

\[
F^{(k)}(\rho)=\sum_{n\geq k}\frac{n!}{(n-k)!}a_n\rho^{n-k}.
\tag{3.6}
\]

Choose a sufficiently small \(h>0\) in the convergence disk of the
Taylor series centered at \(\rho\). Its coefficients are nonnegative by
(3.6). Interchanging nonnegative sums gives

\[
\begin{aligned}
\sum_{k\geq0}\frac{F^{(k)}(\rho)}{k!}h^k
&=\sum_{n\geq0}a_n\sum_{k=0}^n\binom nk\rho^{n-k}h^k\\
&=\sum_{n\geq0}a_n(\rho+h)^n<\infty,
\end{aligned}
\tag{3.7}
\]

contradicting the definition of \(\rho\). \(\square\)

\emph{Proof of Theorem 1.1.} By Lemma 3.1 the series for \(Z\)
represents the actual solution near zero. Denote its convergence radius
by \(\rho\). If \(\rho<R^2\), the actual solution is regular at the
positive point \(\rho\). On an interval around that point, (3.1) is an
ordinary differential equation with analytic right-hand side after
solving for \(Z''\). It therefore has a locally holomorphic solution
with the actual values of \(Z\) and \(Z'\) there. Ordinary uniqueness
identifies this solution with the continuation from zero. This
contradicts Lemma 3.2. Thus \(\rho\geq R^2\).

Conversely, \(z(r)\) tends to infinity as \(r\uparrow R\): by Lemma 2.1
it is increasing, and a finite limit would make
\(w(r)=\int_0^r sz(s)\,ds\) bounded. Therefore \(\rho>R^2\) is
impossible. We have proved \(\rho=R^2\). Integration gives the second
series in (1.3). Dividing its coefficients by a factor linear in the
index leaves the convergence radius unchanged. All derivative assertions
now follow by differentiating convergent positive series on compact
subintervals of \([0,R^2)\). The Hessian assertion was proved after
Lemma 2.1. \(\square\)

The continuation step is essential: positivity of finitely many
coefficients, or even a formal positive recurrence without convergence
control, would not imply a global concavity theorem.

\Needspace{8\baselineskip}\section{4. Convexity of the square root and
sharpness}\label{convexity-of-the-square-root-and-sharpness}

Put \(a_j=c_{j-1}/(2j)>0\). Theorem 1.1 gives

\[
\sqrt{w(r)}=\left\|\bigl(\sqrt{a_j}\,r^j\bigr)_{j\geq1}\right\|_{\ell^2},
\qquad 0\leq r<R.
\tag{4.1}
\]

Every coordinate on the right is nonnegative, nondecreasing and convex
as a function of \(r\). Coordinatewise convexity, monotonicity of the
\(\ell^2\) norm on nonnegative sequences, and the triangle inequality
show that \(r\mapsto\sqrt{w(r)}\) is convex and nondecreasing.
Equivalently one may first prove this for the finite partial sums and
pass to the finite pointwise limit on \([0,R)\).

For strictness it is useful to give an exact identity. At \(r>0\) let

\[
p_j=\frac{a_jr^{2j}}{w(r)},\qquad
m=\sum_{j\geq1}jp_j,\qquad
v=\sum_{j\geq1}(j-m)^2p_j.
\tag{4.2}
\]

These are a probability distribution, its mean and its variance; all
moments are finite inside the convergence interval. Direct
differentiation gives

\[
\frac{r^2\bigl(2ww''-(w')^2\bigr)}{4w^2}
=2v+m(m-1)>0.
\tag{4.3}
\]

The strict inequality holds because \(m>1\), as \(a_2>0\). Thus
\((\sqrt w)''>0\) for \(r>0\). Its tangential Hessian eigenvalue is
\((\sqrt w)'/r>0\).

\emph{Proof of Theorem 1.2.} A convex nondecreasing function of the
Euclidean norm is convex. Applying this fact to (4.1) proves convexity
of \(x\mapsto\sqrt{w(|x|)}\) across the center as well. The Hessian
assertions away from the center follow from (4.3). For \(\beta\geq1/2\),
the function \(s\mapsto s^{2\beta}\) is convex and nondecreasing on
\([0,\infty)\); composing it with \(\sqrt w\) proves the forward
implication of (1.6).

For necessity, the central expansion is

\[
w(r)=\frac A2 r^2+O(r^4).
\tag{4.4}
\]

If \(0<\beta<1/2\), the restriction \(w(r)^\beta\) to a positive radial
segment has strictly negative second derivative for sufficiently small
\(r>0\). Indeed its leading term is \((A/2)^\beta r^{2\beta}\), whose
second derivative has the negative factor \(2\beta(2\beta-1)\); the
differentiated remainder has higher order. It cannot be convex. Finally,
(4.4) gives \(\sqrt{w(|x|)}=\sqrt{A/2}|x|+O(|x|^3)\), which describes
the stated cusp. \(\square\)

The same series gives more than a Hessian bound. For every integer
\(k\geq1\), the radial differential operator
\(\mathcal D=(2r)^{-1}d/dr\) satisfies \(\mathcal D^k w(r)>0\) on
\(0<r<R\), with a positive limit at zero. These infinitely many sign
inequalities refer to differentiation with respect to squared radius,
not arbitrary mixed spatial derivatives.

\Needspace{8\baselineskip}\section{5. Boundary behavior and the optimal positive
power}\label{boundary-behavior-and-the-optimal-positive-power}

Because \(f\) is bounded on the compact range of \(U\), (2.1) gives a
finite limit of \(U'\) at \(R\). The differential equation then gives a
continuous limit of \(U''\) there. The next argument rules out a zero
boundary slope without assuming a finite value of \(f'(0)\).

\Needspace{6\baselineskip}\textbf{Lemma 5.1 (Simple boundary zero).} There is \(b>0\) such that

\[
U(r)=b(R-r)+O((R-r)^2),\qquad
U'(r)=-b+O(R-r),
\tag{5.1}
\]

and

\[
z(r)\sim\frac{1}{R(R-r)}\quad\text{as }r\uparrow R.
\tag{5.2}
\]

\emph{Proof.} Write \(v=-U'>0\). Near \(R\) we have \(0<U<1\), and (2.1)
gives \(U''>0\). Thus \(v\) is decreasing there. Set

\[
F(s)=\frac{s^2}{2}(\log s^2-1),\qquad
E(r)=\frac{v(r)^2}{2}+F(U(r)).
\tag{5.3}
\]

The energy identity is \(E'=-(N-1)v^2/r\). Suppose \(U'(R)=0\). Then
\(E(R)=0\), and for \(r\) sufficiently close to \(R\),

\[
E(r)=\int_r^R\frac{N-1}{s}v(s)^2\,ds
\leq C(R-r)v(r)^2.
\tag{5.4}
\]

Choose this neighborhood so that \(C(R-r)\leq1/4\). Equation (5.4)
implies

\[
\frac{v(r)^2}{4}\leq-F(U(r)),\qquad
v(r)\leq C_1U(r)\sqrt{\log(1/U(r))}.
\tag{5.5}
\]

Since \(U\) is strictly decreasing, it can be used as the variable of
integration on this interval. The remaining time to reach zero would
satisfy

\[
R-r=\int_0^{U(r)}\frac{ds}{v(U^{-1}(s))}
\geq\frac1{C_1}\int_0^{U(r)}\frac{ds}{s\sqrt{\log(1/s)}}
=\infty,
\tag{5.6}
\]

a contradiction. When \(N=1\), the integral in (5.4) is zero and the
same argument applies. Thus \(b=-U'(R)>0\). The bounded second
derivative gives (5.1), and \(z=-U'/(rU)\) gives (5.2). \(\square\)

\emph{Proof of Theorem 1.3.} For \(r>0\),

\[
(U^\alpha)''=\alpha U^\alpha\bigl(\alpha(w')^2-w''\bigr),
\qquad
\frac{(U^\alpha)'}r=-\alpha U^\alpha\frac{w'}r<0.
\tag{5.7}
\]

Thus the radial Hessian condition is precisely \(\alpha\leq q(r)\),
where \(q=w''/(w')^2\). The condition at the center holds automatically,
because \(D^2u^\alpha(0)=-\alpha U(0)^\alpha A I_N\). By Lemma 2.1,
\(q>0\). Equation (4.4) gives \(q(r)\to\infty\) at zero, while Lemma 5.1
gives

\[
q(r)=1-\frac{U(r)U''(r)}{U'(r)^2}\longrightarrow1
\quad\text{as }r\uparrow R.
\tag{5.8}
\]

Moreover \(U''>0\) near \(R\), so \(q<1\) there. Choose any such point.
The limits at both endpoints and continuity imply that the global
infimum is achieved on a compact interior interval; its value is
strictly between zero and one. Equations (5.7) prove (1.7)-\/-(1.8), and
continuity extends concavity to the closed ball. \(\square\)

\Needspace{8\baselineskip}\section{6. Coefficient comparison and two
consequences}\label{coefficient-comparison-and-two-consequences}

The recurrence contains information about variation with the central
height. Fix \(N\) and put \(\delta=A-1>0\). Write \(c_n(\delta)\) for
its coefficients.

\Needspace{6\baselineskip}\textbf{Proposition 6.1 (Positive parameter dependence).} For every
\(n\geq1\), \(c_n(\delta)\) is a polynomial in \(\delta\) with
nonnegative coefficients and zero constant term. If two solutions have
\(0<\delta_1<\delta_2\) and \(\theta=\delta_2/\delta_1\), then on their
common radial interval,

\[
z_2(r)-1\geq\theta\bigl(z_1(r)-1\bigr),
\tag{6.1}
\]

and

\[
w_2(r)-\frac{r^2}{2}
\geq\theta\left(w_1(r)-\frac{r^2}{2}\right).
\tag{6.2}
\]

\emph{Proof.} We have \(c_1=(\delta+\delta^2)/(N+2)\). For \(n\geq1\),
rewrite (1.4) as

\[
c_{n+1}=
\frac{\bigl(2n+1+2(n+1)\delta\bigr)c_n
+(n+1)\displaystyle\sum_{i=1}^{n-1}c_ic_{n-i}}
{(n+1)(2n+N+2)}.
\tag{6.3}
\]

The polynomial assertion follows by induction. Every monomial of
positive degree satisfies \(\delta_2^k\geq\theta\delta_1^k\). Apply this
coefficientwise, including the constant-in-\(r\) term \(c_0-1=\delta\),
and sum the convergent series to obtain (6.1). Integrating
\(r(z_i(r)-1)\) gives (6.2). \(\square\)

\Needspace{6\baselineskip}\textbf{Corollary 6.2 (Known shooting order and uniqueness).} If
\(U_2(0)>U_1(0)\), their first zero radii satisfy \(R_2<R_1\).
Consequently (1.1) has exactly one positive solution for every \(R>0\).

\emph{Proof.} Suppose \(R_2\geq R_1\). By (6.1) and divergence of
\(z_1\) at \(R_1\), strict inequality \(R_2>R_1\) is impossible. If
\(R_2=R_1\), (5.2) gives \(z_2/z_1\to1\), contradicting (6.1) and
\(\theta>1\). This proves the strict shooting order. Two distinct
positive solutions on the same ball are radial and have either different
central values, contradicted by this order, or equal central values,
contradicted by initial-value uniqueness. Existence is the prior input
stated in Section 2. \(\square\)

As emphasized in Section 1.1, this is an alternative proof of a prior
theorem. The coefficientwise parameter comparison (6.1)-\/-(6.2) is the
additional conclusion of the present method.

\Needspace{6\baselineskip}\textbf{Corollary 6.3 (Products of balls).} Let \(u_i\) be positive
solutions of (1.1) on \(B_{R_i}\subset\mathbb R^{N_i}\), for
\(1\leq i\leq k\). Then

\[
u(x_1,\ldots,x_k)=\prod_{i=1}^k u_i(x_i)
\tag{6.4}
\]

solves the same equation on the product domain, and
\(-\sqrt{-\log(u/\|u\|_\infty)}\) is concave there.

\emph{Proof.} Dividing the Laplacian of the product by the product gives
\(-\Delta u/u=\sum_i\log u_i^2=\log u^2\). The negative normalized
logarithm is \(\sum_i w_i\). Its square root is the Euclidean norm of
the vector \((\sqrt{w_i})_{i=1}^k\). Each component is nonnegative and
convex in its corresponding variables by Theorem 1.2. The argument of
(4.1) proves convexity of this norm. \(\square\)

This includes the interval-product construction already treated in
{[}GMS{]}; the factorization itself is elementary and is not a separate
novelty claim.

\Needspace{8\baselineskip}\section{7. Further questions}\label{further-questions}

The argument uses the squared-radius equation and does not extend the
square-root-log conclusion to arbitrary convex domains. Whether that
conclusion survives controlled deformations of a ball, and what
geometric obstructions can prevent it, remain natural questions. The
positivity proof also depends on the particular coefficient structure in
(3.1); a classification of nonlinearities or radial potentials admitting
a comparable positive recurrence would be useful.

Formula (1.8) does not evaluate the optimal positive concavity exponent.
Determining its dependence on the radius in dimensions \(N\geq2\),
especially its sharp limiting behavior, is a further problem. These
questions and the related consequences above belong to the same research
direction; no result about them is asserted here.

\raggedright\Needspace{8\baselineskip}\section{References}\label{references}

{[}BC{]} M. Ben Chrouda, \emph{Uniqueness and Liouville type results for
radial solutions of some classes of k-Hessian equations}, Electronic
Journal of Qualitative Theory of Differential Equations \textbf{2022},
no. 62, 1-\/-17.
\href{https://doi.org/10.14232/ejqtde.2022.1.62}{doi:10.14232/ejqtde.2022.1.62}.

{[}GMS{]} M. Gallo, S. Mosconi and M. Squassina, \emph{Power law
convergence and concavity for the logarithmic Schrödinger equation},
Mathematische Annalen \textbf{395} (2026), article 21, 60 pp.
\href{https://doi.org/10.1007/s00208-026-03452-2}{doi:10.1007/s00208-026-03452-2}.
\href{https://arxiv.org/abs/2411.01614v2}{arXiv:2411.01614v2}.

{[}GMoS{]} M. Gallo, R. Moraschi and M. Squassina, \emph{Quantitative
and exact concavity principles for parabolic and elliptic equations},
\href{https://arxiv.org/abs/2504.09494v2}{arXiv:2504.09494v2}. Also
compared with the
\href{https://www.dmf.unicatt.it/~squassin/papers/lavori/parabconc.pdf}{41-page
author manuscript}, accessed September 15, 2026.

{[}LSZ{]} T. Liu, X. Sun and W. Zou, \emph{Uniqueness of bound states to
the logarithmic Schrödinger equation},
\href{https://arxiv.org/abs/2606.19077v2}{arXiv:2606.19077v2}, June 22,
2026.

\Needspace{8\baselineskip}\section{Preparation note}\label{preparation-note}

This manuscript was prepared with OpenAI Codex assistance in literature
searches, proof development and writing. It is a research preprint. No
independent human verification or journal acceptance is claimed.

\end{document}
